\section{보드 펌웨어 (Pico 기반) --- 코드 분석}

본 장은 \code{Pico/} 디렉터리의 MicroPython 펌웨어를 부팅 흐름, 명령 처리,
하드웨어 추상화, 설정 관리의 네 측면에서 분석한다.

\subsection{부팅 및 메인 루프 (\code{main.py})}

펌웨어 진입점은 \code{AresRover} 클래스이다. 생성자에서 \code{UART(0)}
(9600\,baud, GP0/GP1)와 \code{CommandProcessor}, 하드웨어 싱글턴 \code{robot}을
초기화한다. \code{run()}은 부팅 시퀀스 후 메인 루프를 무한 반복한다.

\begin{lstlisting}[style=py, caption={메인 루프 골격}]
def run(self):
    self.boot()                       # 안전 정지, 부팅음, OLED "ARES READY"
    while self.is_running:
        line = self._read_uart_line() # 개행까지 버퍼링(타임아웃 내)
        if line and not self._is_status_message(line):  # '+' 상태 메시지 필터
            resp = self._process_command(line)
            if self._needs_response(line):
                self._send_response(resp)               # 20바이트 청크 전송
        self.robot.buzzer.update()    # 비차단 부저 자동 정지
        time.sleep_ms(MAIN_LOOP_DELAY_MS)               # 약 10ms 주기
\end{lstlisting}

핵심 동작 특성:

\begin{itemize}
  \item \textbf{다중 청크 명령 조립}: 수신 버퍼(\code{MAX\_BUFFER\_SIZE}, 512바이트)에
        바이트를 누적하고 개행 종단을 폴링한다. 여러 BLE 청크에 걸쳐 도착하는
        긴 명령(\code{BATCH}, LED 패턴, \code{SYS\_SET})을 위해 라인 완성까지
        타임아웃(\code{RECEIVE\_TIMEOUT\_MS}, 약 500\,ms)을 둔다.
  \item \textbf{상태 메시지 필터}: \code{+}로 시작하는 BLE 모듈 상태 메시지를 무시한다.
  \item \textbf{Fire-and-forget}: 응답을 보내지 않는 명령 집합
        (\code{NO\_RESPONSE\_CMDS})을 두어, 웹 프런트엔드와의 명령-응답
        오정렬을 방지한다.
  \item \textbf{비차단 부저}: 음향 명령은 즉시 반환되고, 루프마다 호출되는
        \code{buzzer.update()}가 경과 시간을 검사해 자동 정지시킨다.
\end{itemize}

한 번의 루프 반복에서 명령이 거치는 수명주기는 그림~\ref{fig:fwloop}과 같다.
수신\textbf{·}필터\textbf{·}처리\textbf{·}응답\textbf{·}부저 갱신의 5단계가 약 10\,ms 주기로
반복되며, 응답 여부 판단(\code{\_needs\_response})이 웹과의 명령-응답 정합성을
지키는 분기점이다.

\begin{diagram}{펌웨어 메인 루프 명령 수명주기}{fig:fwloop}
\node[boxw, text width=22mm] (rx) {UART 바이트\\수신};
\node[boxw, right=8mm of rx, text width=22mm] (buf) {버퍼 누적\\개행까지};
\node[boxw, right=8mm of buf, text width=22mm] (filt) {'+' 상태\\메시지 필터};
\node[boxw, right=8mm of filt, text width=24mm] (disp) {Command-\\Processor\\.process()};
\node[boxw, below=9mm of disp, text width=24mm] (need) {\_needs\_\\response?};
\node[boxw, left=10mm of need, text width=24mm] (send) {20B 청크\\응답 전송};
\node[boxw, left=12mm of send, text width=22mm] (buz) {buzzer\\.update()};
\draw[flow] (rx)--(buf); \draw[flow] (buf)--(filt); \draw[flow] (filt)--(disp);
\draw[flow] (disp)--(need);
\draw[flow] (need) -- node[lbl,above]{예} (send);
\draw[flow] (send)--(buz);
\draw[flow] (need.south) to[bend left=30] node[lbl,below]{아니오(fire-and-forget)} (buz.south);
\draw[flow] (buz.north) to[bend left=12] node[lbl,left]{다음 반복} (rx.south west);
\end{diagram}

\subsection{명령 처리기 (\code{process\_data.py})}

\code{CommandProcessor.process(data)}는 수신 문자열을 파싱하여 적절한 핸들러로
분배한다. 명령은 쉼표로 인자를 구분하는 텍스트 프로토콜이며, 성공 시 \code{1},
미인식\textbf{·}오류 시 \code{0} 또는 \code{ERROR}를 반환한다.

\subsubsection{디스패치 구조}

\code{process()}는 긴 \code{if/elif} 사슬로, 두 가지 매칭 방식을 혼용한다.
인자 없는 명령은 \emph{완전 일치}(\code{data == "PING"})로, 인자 있는 명령은
\emph{접두 일치}(\code{data.startswith("SYS\_SET")})로 판별한다.

\begin{lstlisting}[style=py, caption={process() 디스패치 사슬(발췌)}]
if   data == "PING":            return self._handle_ping()
elif data == "STOP_ALL":        return self._handle_stop_all()
elif data.startswith("SYS_SET"):return self._handle_sys_set(data)
elif data.startswith("BATCH;"): return self._handle_batch(data)
elif data.startswith("SERVO_tFORWARD"): return self._handle_timed_wheel(data,"forward")
elif data.startswith("tFORWARD"):        return self._handle_timed_wheel(data,"forward")
...
else: return 0   # 미인식 명령
\end{lstlisting}

\begin{infobox}[title=접두 일치의 순서 의존성]
접두 일치는 \textbf{선언 순서가 중요}하다. 더 구체적인 접두(\code{SERVO\_tFORWARD})를
더 일반적인 접두(\code{tFORWARD})보다 먼저 검사해야 한다. 또한
\code{startswith("SING")}처럼 짧은 접두는 다른 명령과의 우연한 일치를
피하도록 배치된다. 새 명령을 추가할 때는 이 순서 규칙을 지켜야 한다.
\end{infobox}

\subsubsection{시스템 \textbf{·} 연결 명령}

\begin{longtable}{L{3.4cm} L{11.0cm}}
\toprule
\rowcolor{tablehead}\textbf{명령} & \textbf{동작} \\
\midrule
\endhead
\code{PING} & 연결 확인. OLED에 ``CONNECTED!''와 장치명 표시, \code{1} 반환 \\
\code{STOP\_ALL} & 비상 정지. 모든 모터/출력 정지, OLED ``EMERGENCY STOP'' \\
\code{GET\_STATUS} & \code{STATUS,거리,자기,온도,여유메모리,업타임} 반환 \\
\code{GET\_SYS} & \code{SYS\_VALUES,속도,충돌거리,자동정지,이름,좌보정,우보정} 반환 \\
\code{GET\_MODULES} & \code{MODULES,모듈:ON|OFF,...} 형식으로 전체 모듈 상태 반환 \\
\code{SYS\_SET,속도,거리,정지,이름} & 시스템 설정 갱신 후 저장 \\
\code{SET\_PIN,이름,번호} & 핀 할당 변경 \\
\code{SET\_MODULE,이름,0|1} & 모듈 활성/비활성 \\
\code{BATCH;cmd1|cmd2|...} & 여러 명령 순차 실행(단일 응답) \\
\code{CALIB\_START} / \code{CALIB\_SET,좌,우} & 휠 직진 보정 시작/적용 \\
\code{SLEEP,초} & 지정 시간 대기(차단) \\
\bottomrule
\end{longtable}

\subsubsection{구동 명령}

\begin{longtable}{L{4.6cm} L{9.8cm}}
\toprule
\rowcolor{tablehead}\textbf{명령} & \textbf{동작} \\
\midrule
\endhead
\code{SERVO\_FORWARD} / \code{FORWARD} 등 & 서보 휠 연속 전진/후진/좌/우회전 \\
\code{SERVO\_STOP} / \code{STOP} & 서보 휠 정지 \\
\code{tFORWARD,초} 등 & 서보 휠 시간 지정 이동(차단), 완료 후 응답 \\
\code{DC\_FORWARD} / \code{MAIN\_FORWARD} 등 & DC 모터 연속 구동 \\
\code{DC\_STOP} & DC 모터 정지 \\
\code{DC\_tFORWARD,초} 등 & DC 모터 시간 지정 구동(차단) \\
\bottomrule
\end{longtable}

서보 휠은 \code{KSWheel}이 좌우 서보 각도를 제어하여 전진/후진/회전을
구현하며, 보정 계수를 곱해 직진성을 교정한다. DC 모터는 \code{max\_speed}
설정값을 기준으로 PWM 듀티를 산출한다.

\subsubsection{출력 \textbf{·} 센서 명령}

\begin{longtable}{L{4.6cm} L{9.8cm}}
\toprule
\rowcolor{tablehead}\textbf{명령} & \textbf{동작} \\
\midrule
\endhead
\code{LED\_ON,번호,밝기} & 단일 LED 켜기(밝기 0.0\textasciitilde 1.0) \\
\code{LED\_OFF,번호|ALL} & 단일/전체 LED 끄기 \\
\code{[v0 v1 v2 v3 v4 v5]} & 6개 LED 밝기 일괄 설정(공백 구분) \\
\code{LED\_PATTERN} & 스와이프 애니메이션 \\
\code{BUZZER\_ON,주파수,지속} & 비차단 톤 시작 \\
\code{SING} & 내장 멜로디 재생 \\
\code{MSG,텍스트} / \code{MSG\_XY,x,y,텍스트} & OLED 텍스트 출력 \\
\code{ICON,이름,x,y} & 아이콘 표시(rover/mars/open\_eye/closed\_eye) \\
\code{CLEAR\_DISPLAY} / \code{CLEAR\_RECT,x,y,w,h} & 화면/영역 지우기 \\
\code{GUN\_FIRE} & BB탄 단발 발사 \\
\code{DISTANCE} & \code{DIST:값} 반환(cm) \\
\code{MAGNET} & \code{MAG:0|1} 반환 \\
\bottomrule
\end{longtable}

\subsection{하드웨어 추상화 (\code{hardware.py})}

\code{RobotHardware}는 싱글턴 패턴(\code{\_instance})으로 구현되어 전역
\code{robot} 인스턴스로 접근된다. 각 모듈 속성은 초기값 \code{None}이며,
\code{system\_config.is\_module\_enabled()} 검사 후 try-except로 초기화한다.
DC 모터는 안전을 위해 가장 먼저 초기화하고 즉시 정지시킨다.

주요 메서드:
\begin{itemize}
  \item \code{get\_temperature()} --- ADC4 기반 섭씨 온도 환산
  \item \code{get\_status\_dict()} --- 거리\textbf{·}자기\textbf{·}온도\textbf{·}여유메모리(\code{gc.mem\_free()})\textbf{·}업타임 딕셔너리
  \item \code{stop\_all()} --- 모듈별 예외 처리를 동반한 전체 안전 정지
\end{itemize}

\subsection{설정 관리 (\code{system\_config.py})}

\code{SystemConfig}(싱글턴)는 \code{DEFAULT\_CONFIG}를 기준값으로 두고,
\code{system.json}이 존재하면 병합하여 덮어쓴다. 구형 \code{calibration.json}이
있으면 보정값을 이관한다.

\begin{lstlisting}[style=json, caption={system.json 주요 항목}]
{
  "max_speed": 80, "collision_dist": 10, "auto_stop": 1,
  "device_name": "ARES",
  "left_calibration": 100, "right_calibration": 100,
  "enable_wheel": 1, "enable_dcmotor": 1, "enable_buzzer": 1,
  "enable_distance": 1, "enable_magsensor": 1, "enable_leds": 1,
  "enable_gun": 1, "enable_oled": 1,
  "pin_wheel_left": 13, "pin_wheel_right": 12,
  "pin_dcmotor_dir": 8, "pin_dcmotor_pwm": 9, "pin_buzzer": 15,
  "pin_distance_trig": 6, "pin_distance_echo": 7,
  "pin_magsensor": 26, "pin_gun": 22,
  "pin_leds": "21,20,19,18,17,16",
  "pin_i2c_sda": 4, "pin_i2c_scl": 5
}
\end{lstlisting}

핵심 메서드: \code{is\_module\_enabled()}, \code{enable\_module()},
\code{set\_pin()}(0\textasciitilde 28 검증), \code{save\_config()},
\code{save\_calibration()}, \code{save\_all()}.

\subsection{펌웨어 설계상의 제약과 패턴}

\begin{infobox}[title=주요 제약]
\begin{itemize}
  \item \textbf{RAM 제한}: 대용량 버퍼 회피, 수신 버퍼 512바이트 상한,
        \code{json.indent} 미지원.
  \item \textbf{BLE 20바이트 청크}: 송수신 모두 20바이트 단위.
  \item \textbf{싱글턴 강제}: \code{RobotHardware} 인스턴스는 단 하나.
  \item \textbf{차단형 시간 명령}: \code{tFORWARD} 등 시간 지정 명령은 메인 루프를
        점유하므로, 그 동안 다른 명령\textbf{·}비상정지 응답이 지연될 수 있다.
\end{itemize}
\end{infobox}

\code{Pico/backup/} 및 날짜 표기 백업 파일(예: \code{main\_backup\_*.py})은
레거시\textbf{·}참조용으로, 수정 대상이 아니다.

\begin{teachbox}{펌웨어의 단순함과 그 대가}
\teachhead{설계 의도}
\begin{itemize}
  \item 단일 수신 루프 + \code{if/elif} 디스패치로 \emph{읽기 쉬움}을 택했다.
  \item 부저는 비차단, 시간 명령은 차단으로 처리한 \emph{트레이드오프}가 있다.
  \item 응답 없는 명령(fire-and-forget)으로 명령--응답 어긋남을 막는다.
\end{itemize}
\teachhead{분석할 때 핵심 질문}
\begin{itemize}
  \item 시간 명령(\code{tFORWARD,2})이 도는 2초 동안 비상정지가 왜 늦어질 수 있는가?
  \item 접두 일치 순서를 바꾸면 어떤 명령이 잘못 잡힐까?
\end{itemize}
\teachhead{점검 요소}
\begin{itemize}
  \item \code{process\_data.py}에 새 명령 하나를 추가하되, 더 구체적인 접두를 먼저 두는 순서 규칙을 지켜라.
\end{itemize}
\end{teachbox}

